\chapter{Remaining FYP Plan}
\section{SABLE Continuation: Proposed Research Work}
The following items describe future SABLE research work, not completed mid-stage features or an approved team schedule. Confirm priorities, owners, dates, and the final evaluation protocol with the supervisor and team.
\begin{enumerate}
    \item Freeze the bounded obligation and scenario labels, then validate the ground truth independently.
    \item Implement and document the strongest agreed comparison baselines; compare B0--B4 and relevant existing tools under a reproducible protocol.
    \item Measure false-safe and false-regression behavior, correspondence outcomes, uncertainty quality, and evidence completeness on a larger controlled corpus.
    \item Extend Terraform/IAM coverage only where justified by the confirmed research question, and preserve explicit unsupported boundaries.
    \item Run broader integrated workflow validation and report any other module's unresolved status separately from SABLE results.
\end{enumerate}
These experiments have not been completed by the focused tests reported in Chapter 8. No performance, novelty, or external-validity conclusion is claimed.

\section{Remaining Scope}
The mid-evaluation implementation is a bounded prototype. SABLE does not cover all Terraform expressions or AWS IAM semantics and does not evaluate deployed cloud state. The team should confirm module-specific remaining scope from each owner's implementation and evidence before setting final-project acceptance criteria.
\section{Remaining Modules and Features}
The verified integration branch combines CAVR, SATRA-RV, SABLE, and the shared workflow. Remaining module work and ownership must be recorded from the team's accepted plans; this report does not infer individual tasks from module presence alone.
\section{Remaining Iterations}
FYP-2 is proposed to focus on the agreed comparative evaluation, stronger independent validation of scenario labels, and additional integration checks. These are future activities, not completed results.
\section{Implementation Plan}
After supervisor and team agreement, record owner-assigned work items for each approved research question, maintain the bounded security models, and review changes against the corresponding requirements and regression tests. No dates or owners are asserted here because an approved schedule was not present in the reviewed records.
\section{Testing and Validation Plan}
\begin{sloppypar}
Retain focused module regressions, Docker-enabled sandbox checks, frontend build/render checks, and the combined suite on the same integration commit. Independently review the benchmark labels and add cases only when their expected result follows from the stated obligation and supported policy model. Report skips and environment limitations alongside every result.
\end{sloppypar}
\section{Deployment Plan}
The current demonstration runs locally with supplied fixtures and does not apply Terraform or contact AWS. Any later deployment or live-cloud evaluation requires a separately approved test environment, credentials policy, and scope; none is claimed for this mid-evaluation report.
\section{Tentative Timeline}
The reviewed repository did not contain an approved FYP-2 calendar or dated owner plan. The team should populate this section from the supervisor-approved schedule rather than infer dates from the current implementation history.
\section{Expected Final Deliverables}
Expected final-project outputs, subject to supervisor approval, are the integrated implementation, traceable requirements/design and test evidence, a reproducible comparative evaluation, and the final report and demonstration. This list is a proposal for planning, not evidence that those deliverables are complete.
